\chapter*{Introduction Générale}
\addcontentsline{toc}{chapter}{Introduction Générale}

\noindent
Le stage constitue une étape importante dans le parcours d'un étudiant en ingénierie logicielle. Il permet de confronter les connaissances acquises durant la formation aux besoins réels d'un projet professionnel, tout en développant des compétences techniques, organisationnelles et relationnelles.

\noindent
Dans le cadre de ce stage, le projet confié porte sur la conception et la réalisation de \projectName{}, une plateforme web de recrutement dédiée au secteur CHR au Maroc. Ce secteur regroupe les cafés, hôtels et restaurants, qui rencontrent souvent des besoins rapides en personnel et s'appuient encore fréquemment sur des méthodes de recrutement informelles : messages WhatsApp, recommandations, réseaux sociaux ou CV papier.

\noindent
\projectName{} vise à proposer une solution numérique centralisée permettant de mieux organiser le recrutement. La plateforme offre aux candidats un espace pour compléter leur profil, ajouter leur CV, consulter les offres et suivre leurs candidatures. Elle offre aux recruteurs un espace pour publier des offres, ajouter des quiz métier, consulter les profils reçus et gérer les statuts des candidatures. Enfin, un espace administrateur permet la modération des offres et le suivi global de l'activité.

\noindent
Le présent rapport décrit le contexte du stage, l'analyse des besoins, les choix méthodologiques et techniques, la conception du système ainsi que les principales réalisations effectuées. Il met également en évidence les difficultés rencontrées, les solutions apportées et les perspectives d'amélioration de la plateforme.
